% Part: computability
% Chapter: recursive-functions
% Section: trees

\documentclass[../../../include/open-logic-section]{subfiles}

\begin{document}

\olfileid{cmp}{rec}{tre}
\olsection{વૃક્ષો}

ક્યારેક વૃક્ષોને પ્રાકૃતિક સંખ્યાઓ તરીકે રજૂ કરવું ઉપયોગી છે. જેમ
અનુક્રમોને સંખ્યાઓથી રજૂ કરીએ છીએ અને તેમનાં ગુણધર્મો તથા ક્રિયાઓને
તેમના કોડ પરના આદિમ પુનરાવર્તી સંબંધો અને વિધેયોથી રજૂ કરીએ છીએ,
તેમ વૃક્ષો માટે પણ કરી શકીએ. આ માટે અનુક્રમો અને તેમના કોડ
વાપરીશું. વૃક્ષમાં માત્ર એક જ ગાંઠ હોઈ શકે છે (જેને નામચિહ્ન હોય
કે ન હોય), અથવા કેટલીક ઉપવૃક્ષો સાથે જોડાયેલી ગાંઠ હોઈ શકે છે
(જેને નામચિહ્ન હોય કે ન હોય). એ ગાંઠને વૃક્ષનું \emph{મૂળ} અને
તેની સાથે જોડાયેલાં ઉપવૃક્ષોને તેનાં \emph{તત્કાલ ઉપવૃક્ષો} કહીએ.

વૃક્ષને પુનરાવર્તી રીતે અનુક્રમ $\tuple{k, d_1, \dots, d_k}$ વડે
સંકેતીકૃત કરીએ છીએ. અહીં $k$ એ તત્કાલ ઉપવૃક્ષોની સંખ્યા છે અને
$d_1$, \dots,~$d_k$ તેમના કોડ છે. ગાંઠોને નામચિહ્નો હોય, તો
તેમને તત્કાલ ઉપવૃક્ષોની પાછળ ઉમેરી શકાય. તેથી માત્ર એક ગાંઠનું,
નામચિહ્ન~$l$ ધરાવતું વૃક્ષ $\tuple{0,l}$થી સંકેતીકૃત થશે.
નામચિહ્ન~$l_1$ ધરાવતા મૂળને નામચિહ્નો $l_2$, $l_3$ ધરાવતી બે
એકગાંઠવાળી ઉપવૃક્ષો જોડાયેલી હોય, તો તેનો કોડ
$\tuple{2,
  \tuple{0, l_2}, \tuple{0, l_3}, l_1}$ થશે.

\begin{prop}
  \ollabel{prop:subtreeseq}
  વિધેય $\fn{SubtreeSeq}(t)$ આદિમ પુનરાવર્તી છે. તે એવા
  અનુક્રમનો કોડ આપે છે જેના ઘટકો કોડ~$t$ ધરાવતા વૃક્ષનાં બધાં
  ઉપવૃક્ષોના કોડ છે.
\end{prop}

\begin{proof}
  પહેલાં નોંધો કે $\fn{ISubtrees}(t) = \fn{subseq}(t, 1, (t)_0)$
  આદિમ પુનરાવર્તી છે અને વૃક્ષ~$t$નાં તત્કાલ ઉપવૃક્ષોના કોડ આપે
  છે. હવે સહાયક વિધેય $\fn{hSubtreeSeq}(t,n)$ વ્યાખ્યાયિત કરી
  શકીએ, જે મૂળથી વધુમાં વધુ $n$ કડીઓ દૂરનાં બધાં ઉપવૃક્ષોનો
  અનુક્રમ ગણે છે. મૂળથી $0$ કડી દૂરનાં~$t$નાં ઉપવૃક્ષોનો
  અનુક્રમ---અર્થાત્~$t$ના મૂળથી શરૂ થતો અનુક્રમ---માત્ર~$t$ ધરાવે છે.
  વૃક્ષ~$t$ના મૂળથી વધુમાં વધુ $n+1$ કડીઓ દૂરનાં બધાં ઉપવૃક્ષોનો
  અનુક્રમ મેળવવા, વધુમાં વધુ $n$ કડીઓ દૂરનાં ઉપવૃક્ષોના અનુક્રમ
  સાથે આ $n$ સ્તર સુધીનાં ઉપવૃક્ષોનાં બધાં તત્કાલ ઉપવૃક્ષોનો
  અનુક્રમ જોડીએ. આ બધાની યાદી બનાવવા નોંધો કે
  $f(x)$ એ અનુક્રમોના કોડ આપતું આદિમ પુનરાવર્તી વિધેય હોય, તો
  પહેલાં $k$ ઘટકો પરનું જોડાણ $g_f(s, k) = f((s)_0) \concat
  \dots \concat f((s)_{k-1})$ પણ આદિમ પુનરાવર્તી છે; શૂન્ય
  ઘટકો માટે તે ખાલી અનુક્રમ છે:
    \begin{align*}
      g_f(s, 0) & = \emptyseq\\
      g_f(s, k+1) & = g_f(s, k) \concat f((s)_k)
    \end{align*}
    \sourcecorrection{OLCMP-008}{સ્થિર અંગ્રેજી મૂળ
    $g_f(s,k)$માં $0$થી $k$ સુધી, એટલે $k+1$ ઘટકો જોડે છે અને
    સમીકરણોમાં $g$ લખે છે; પછી $h(s)$ માટે $\len{s}$ વાપરે છે.
    ગુજરાતી વ્યાખ્યામાં પહેલાં $k$ ઘટકો જોડીએ છીએ, $k=0$ માટે
    ખાલી અનુક્રમ રાખીએ છીએ, અને સમીકરણોમાં $g_f$ રાખીએ છીએ.
    તેથી પછીનો $g_{\fn{ISubtrees}}(s,\len{s})$ ખરેખર $s$ના
    બધા ઘટકો પર જ ચાલે છે.}
    દાખલા તરીકે, $s$ એ વૃક્ષોનો અનુક્રમ હોય, તો
    $h(s) = g_{\fn{ISubtrees}}(s, \len{s})$ એ~$s$ના ઘટકોનાં
    તત્કાલ ઉપવૃક્ષોનો અનુક્રમ આપે છે. આ વડે $\fn{hSubtreeSeq}$ને
    આમ વ્યાખ્યાયિત કરી શકીએ:
    \begin{align*}
      \fn{hSubtreeSeq}(t, 0) & = \tuple{t} \\
      \fn{hSubtreeSeq}(t, n+1) & = \fn{hSubtreeSeq}(t, n) \concat
      h(\fn{hSubtreeSeq}(t, n)).
    \end{align*}
    \sourcecorrection{OLCMP-009}{સ્થિર અંગ્રેજી મૂળ આ સહાયક
    વિધેયને મૂળથી બરાબર $n$ કડીઓ દૂરનાં ઉપવૃક્ષોનો અનુક્રમ કહે છે.
    પ્રદર્શિત પુનરાવર્તી સમીકરણ અગાઉની સમગ્ર યાદી સાથે નવા
    ઉપવૃક્ષો જોડે છે; તેથી તેમાં વધુમાં વધુ $n$ કડીઓ દૂરનાં
    ઉપવૃક્ષો રહે છે. ગુજરાતી વર્ણનમાં આ સંચિત અર્થ રાખ્યો છે.}
    કોડ~$t$ ધરાવતા વૃક્ષમાં ઉપવૃક્ષોનું મહત્તમ સ્તર, એટલે મૂળથી
    છેડાની ગાંઠ સુધીનું મહત્તમ અંતર, કોડ~$t$થી ઉપરથી સીમિત છે.
    તેથી કોડ~$t$ ધરાવતા વૃક્ષનાં બધાં ઉપવૃક્ષોના કોડનો અનુક્રમ
    $\fn{hSubtreeSeq}(t, t)$ આપે છે.
\end{proof}

\begin{prob}
  \olref[cmp][rec][tre]{prop:subtreeseq}ની સાબિતીમાં
  $\fn{hSubtreeSeq}$ની વ્યાખ્યા સામાન્ય રીતે પુનરાવૃત્તિઓ
  ધરાવે છે. મળતી યાદીમાં કોઈ ઉપવૃક્ષનો કોડ માત્ર એક વાર જ આવે,
  તેની ખાતરી આપતી બીજી વ્યાખ્યા આપો.
\end{prob}

\end{document}
